\section{A uniform family of double-tilt aisles}
\label{sec:aisles}

Our goal is to construct bounded aisles $U_\eta$ for which
$\Phi U_\eta\subset U_\theta$ holds uniformly with independent parameters.
We begin with Toda's proposed two tilts for the quintic Gepner charge
\cite[Sections~3.4--3.6]{TodaQuintic2013}, and vary the second cut in two
numerical directions. The slope estimate below supplies the strict margins
that make the comparison survive these variations.
The two first-tilt slopes $-1/2$ and $1/2$ reflect the order of the factors
in $\Phi$: tensoring by $\OO_X(1)$ moves the first cut from $-1/2$ to
$1/2$, and the spherical twist will take the resulting aisle into one
constructed at $-1/2$.

We first isolate the geometric input.  Slope always means
$\mu(F)=c(F)/r(F)$ for a torsion-free sheaf of positive rank.  We use
ordinary slope Harder--Narasimhan filtrations, writing $\mu^+(F)$ and
$\mu^-(F)$ for their largest and smallest slopes.  For an arbitrary
sheaf, its torsion subsheaf is removed before these slopes are taken.

\begin{lemma}[Slope input]\label{lem:slope-input}
Put
\[
 f(x)=\max\left\{-\frac{3x}{25},\ \frac{x}{2}-\frac7{25},\
                    \frac{28x}{25}-\frac{31}{50}\right\}
                    \qquad(0\leq x\leq1).
\]
Every torsion-free slope-semistable sheaf $F$ on $X$ with
$|\mu(F)|\leq1$ satisfies
\begin{equation}\label{eq:slope-envelope}
                 \frac{d(F)}{r(F)}\leq f(|\mu(F)|).
\end{equation}
The estimate extends under integral twists by
$(x,y)\mapsto(x+n,y+nx+n^2/2)$.  In particular,
$f(x)\leq0$ for $0\leq x\leq1/2$, and $f(1/2)=-3/100$.
\end{lemma}

\begin{proof}
We use Xu's Theorem~1.3 \cite{Xu2026}.  Its hypotheses are precisely a
smooth complex quintic threefold with its hyperplane polarization and a
torsion-free slope-semistable sheaf of slope in $[-1,1]$; its coordinates
$\mu_H$ and $\xi_H$ are our $c/r$ and $d/r$.  On $[0,1/2]$ its upper
bound $g$ is
\[
g(x)=
\begin{cases}
-x/2,&0\leq x\leq1/4,\\
x/2-1/4,&1/4\leq x\leq5/13,\\
-3x/20,&5/13\leq x\leq6/13,\\
x/2-3/10,&6/13\leq x\leq1/2.
\end{cases}
\]
The first and third lines are at most $-3x/25$; the second is also at
most $-3x/25$ because $5/13<25/62$; and the last is at most
$x/2-7/25$.  Thus $g\leq f$ on this half interval.  The full upper
bound $g$ in Xu's theorem and the maximum $f$ both satisfy
\[
 h(1-x)=h(x)+\frac12-x.
\]
For $f$, this reflection exchanges the first and third affine pieces
and fixes the second.  This proves the comparison on $[0,1]$, and the
theorem uses $|\mu|$ on the negative strip.

For completeness, the symmetry is also compatible with the sheaf
operations used here.  Twisting by $\OO_X(n)$ has the displayed effect
on $(x,y)$ and preserves slope semistability.  If $F$ is not reflexive,
passing to $F^{**}$ preserves slope semistability and $r,c$, and only
increases $d$.  For a reflexive sheaf on a smooth threefold the higher
sheaf Ext terms of its derived dual are supported in dimension zero,
so its ordinary dual has characters $(r,-c,d)$ through degree two and
is slope-semistable.  The operations of dualizing and twisting by
$\OO_X(1)$ therefore give $(x,y)\mapsto(1-x,y+1/2-x)$ without changing
the validity of an upper bound.  Finally the three defining lines are
nonpositive on $[0,1/2]$, and their maximum at $1/2$ is $-3/100$.
\end{proof}

We use aisles in the convention $\D^{\leq0}$: they are closed under
$[1]$, and the heart of an aisle $V$ is $V\cap(V[1])^\perp$, where
$\perp$ denotes right Hom orthogonality.  Recall explicitly the tilt
construction of Happel--Reiten--Smal\o\ \cite{HappelReitenSmalo1996}.
A torsion pair $(\mathcal T,\mathcal F)$ in the heart of
a bounded aisle $V$ gives the aisle
\begin{equation}\label{eq:aisle-tilt}
 V^\sharp=\{E\in V:H^0_V(E)\in\mathcal T\},\qquad
 \mathcal H^\sharp=\langle\mathcal F[1],\mathcal T\rangle,
 \qquad V[1]\subset V^\sharp\subset V.
\end{equation}
Here and below angle brackets denote extension closure.  To construct
the new truncation triangle, first truncate at zero for $V$, then take
the inverse image of the torsion subobject of its zeroth cohomology.
The remaining triangle has terms from $\mathcal F$ and the old positive
degrees.  The old orthogonality and
$\Hom(\mathcal T,\mathcal F)=0$ give the required new orthogonality.
Shift closure and the two inclusions prove that this is a bounded
$t$-structure.  The same truncations show conversely that an aisle
between $V[1]$ and $V$ gives a torsion pair in the heart of $V$.

For $b\in\{-1/2,1/2\}$, let $\mathcal C_{\leq b}$ consist of zero
and the torsion-free sheaves with $\mu^+\leq b$, and let
$\mathcal C_{>b}$ consist of sheaves whose torsion-free quotient has
$\mu^->b$, together with all torsion sheaves.  Ordinary slope
filtrations give the torsion pair
$(\mathcal C_{>b},\mathcal C_{\leq b})$ in $\Coh(X)$.  Define the
first tilted heart and its aisle by
\[
 \mathcal B_b=\langle\mathcal C_{\leq b}[1],\mathcal C_{>b}\rangle,
 \qquad W_b=\D^{\leq0}_{\mathcal B_b},
 \qquad p_b=c-br.
\]
On $\mathcal B_b$, $p_b$ takes values in $\frac12\Z_{\geq0}$, by
Lemma~\ref{lem:numerical}.  For $E\in\mathcal B_b$, the equality
$p_b(E)=0$ holds exactly when $H^{-1}(E)$ is slope-semistable of slope
$b$ (or zero), and $H^0(E)$ has dimension at most one.  This follows
by applying $p_b$ to the slope factors of the two cohomology sheaves:
each contribution is nonnegative, and a nonzero divisorial torsion
sheaf has positive $c$.

\begin{lemma}\label{lem:first-tilt-noetherian}
The abelian categories $\mathcal B_{-1/2}$ and $\mathcal B_{1/2}$ are
noetherian.
\end{lemma}

\begin{proof}
We give the rational first-tilt argument; compare
\cite[proof of Lemma~3.2.4]{BayerMacriToda2014}.
Consider epimorphisms $E_0\twoheadrightarrow E_1\twoheadrightarrow
E_2\twoheadrightarrow\cdots$ in $\mathcal B_b$, with kernels $K_i$
for the individual arrows.  The nonnegative discrete numbers
$p_b(E_i)$ eventually stabilize, so $p_b(K_i)=0$.  The ordinary
cohomology sheaves $H^0(E_i)$ form a sequence of epimorphisms and
eventually stabilize, since $\Coh(X)$ is noetherian.  The preceding
description of $K_i$ gives $r(K_i)\leq0$, hence the integers $r(E_i)$
are nondecreasing.  They are bounded above by the now fixed rank of
$H^0(E_i)$, so eventually $r(K_i)=0$.  Each remaining $K_i$ is then an
ordinary sheaf of dimension at most one.

The cohomology sequence now reads
\[
 0\longrightarrow H^{-1}(E_i)\longrightarrow H^{-1}(E_{i+1})
   \longrightarrow K_i\longrightarrow0.
\]
These torsion-free sheaves agree in codimension one.  Their injections
extend to isomorphisms of their reflexive hulls, so they form an
increasing chain of coherent subsheaves of a fixed reflexive sheaf.
That chain stabilizes, forcing $K_i=0$.  Thus every such sequence of
epimorphisms stabilizes, as required.
\end{proof}

The first tilted hearts are now available. We define the second cut and choose
constants for three uses: negativity at the slope boundary, exclusion of a
small neighboring slope strip, and a strict comparison between the two cuts.
At zero perturbation and $b=-1/2$, the cut is the imaginary part of the
normalized Gepner charge, as in Toda's candidate heart
\cite[Section~3.6]{TodaQuintic2013}. Fix
\begin{equation}\label{eq:aisle-parameters}
 \begin{gathered}
 \delta=10^{-5},\qquad s_0=\frac{553}{1000},\qquad
 B_\delta=\{\eta=(\eta_0,\eta_1)\in\R^2:|\eta_j|<\delta\},\\
 N_b^\eta=d-bc+ur+\eta_1c+\eta_0r.
 \end{gathered}
\end{equation}
The following exact margins will be useful:
\begin{align}
 m_b&=\frac7{25}-u-\frac32\delta>0,\label{eq:boundary-margin}\\
 m_s&=\frac{13857}{50000}-u-(1+s_0)\delta>0,
                                                    \label{eq:strip-margin}\\
 m_c&=\frac12-\frac{u}{s_0}-3\delta>0.\label{eq:comparison-margin}
\end{align}
For example, $\sqrt5>223606/100000$ implies
$u<138197/500000$, and gives respectively the positive lower bounds
$3591/10^6$, $73047/10^8$, and $8941/55300000$ for these three
expressions.  On a slope-semistable sheaf $F$ of slope $b$,
Lemma~\ref{lem:slope-input} gives
\begin{equation}\label{eq:boundary-negative}
 \frac{N_b^\eta(F)}{r(F)}
 \leq -\frac3{100}-\frac14+u+\frac32\delta=-m_b<0.
\end{equation}
Consequently $N_b^\eta\geq0$ on the $p_b=0$ objects of $\mathcal B_b$;
equality holds precisely for zero-dimensional sheaves, including zero.
Indeed their degree-minus-one part contributes strictly positively
unless it vanishes, while a sheaf of dimension at most one has
$N_b^\eta=d\geq0$, with equality exactly in dimension zero.

Call a nonzero object of $\mathcal B_b$ \emph{high} if $p_b=0$ or
$N_b^\eta>0$, and \emph{low} otherwise.  Thus low means $p_b>0$ and
$N_b^\eta\leq0$.  Set
\[
\begin{split}
 \mathcal T_b^\eta&=\{E:\text{every nonzero quotient of $E$ in
                                  $\mathcal B_b$ is high}\},\\
 \mathcal F_b^\eta&=\{E:\text{every nonzero subobject of $E$ in
                                  $\mathcal B_b$ is low}\}.
\end{split}
\]

\begin{lemma}[The second tilt]\label{lem:second-tilt}
For every $b=\pm1/2$ and $\eta\in B_\delta$, the pair
$(\mathcal T_b^\eta,\mathcal F_b^\eta)$ is a torsion pair.  Its tilt
has bounded heart and aisle
\[
 \mathcal A_b^\eta=\langle\mathcal F_b^\eta[1],\mathcal T_b^\eta\rangle,
 \qquad \mathcal U_b^\eta
   =\{E\in W_b:H^0_{\mathcal B_b}(E)\in\mathcal T_b^\eta\}.
\]
On $\mathcal A_b^\eta$ one has $N_b^\eta\geq0$, and no nonzero object
of this heart has zero numerical class.
\end{lemma}

\begin{proof}
High objects are closed under extensions: if one end term has positive
$N_b^\eta$, so does the extension, and otherwise both have $p_b=0$.
It follows, by taking the induced filtration on a quotient, that
$\mathcal T_b^\eta$ is closed under quotients and extensions.

If an object has a high subobject, choose such a nonzero subobject $A$
with $p_b(A)$ minimal.  When $p_b(A)=0$, all its quotients have
$p_b=0$.  Otherwise $N_b^\eta(A)>0$, and a low quotient $Q$ would
have $p_b(Q)>0$ and $N_b^\eta(Q)\leq0$.  Its kernel would then be a
high subobject with strictly smaller $p_b$, a contradiction.  Thus
$A\in\mathcal T_b^\eta$ in either case.

By Lemma~\ref{lem:first-tilt-noetherian}, every object has a maximal
$\mathcal T_b^\eta$ subobject; sums of such subobjects again belong to
$\mathcal T_b^\eta$, so maximality has its usual torsion meaning.  The
quotient has no high subobject, since such a subobject would contain
a nonzero member of $\mathcal T_b^\eta$ by the preceding paragraph,
and its inverse image would enlarge the maximal torsion subobject.
The quotient therefore lies in $\mathcal F_b^\eta$.  Finally a
nonzero image of a map from $\mathcal T_b^\eta$ to
$\mathcal F_b^\eta$ would be both high and low.  This proves the
torsion pair, and \eqref{eq:aisle-tilt} gives boundedness.

Write an object of the new heart as an extension of $T$ by $F[1]$,
where $T\in\mathcal T_b^\eta$ and $F\in\mathcal F_b^\eta$.  Then
$N_b^\eta(T)\geq0$ and $N_b^\eta(F)\leq0$, proving nonnegativity.
If $N_b^\eta(E)=0$, both terms have $N_b^\eta=0$.  Since $T$ itself
is high when nonzero, this forces $p_b(T)=0$, hence $T$ is
zero-dimensional.  If $F\ne0$, then $p_b(F)>0$ and
$p_b(E)=-p_b(F)<0$.  Thus $[E]=0$ would force $F=0$; it would then
force $T=0$, since $e(T)=\operatorname{length}(T)/5$.  This proves
the last assertion.
\end{proof}

We next compare the two slope cuts with independent parameters.
Subscripts $-$ and $+$ mean $b=-1/2$ and $b=1/2$, respectively.

\begin{lemma}\label{lem:compare-slope-cuts}
For every $\alpha,\theta\in B_\delta$,
\[
                       \mathcal U_+^\alpha\subset\mathcal U_-^\theta.
\]
\end{lemma}

\begin{proof}
Let $\mathcal M$ be the ordinary torsion-free sheaves all of whose
slopes lie in $(-1/2,1/2]$, and let
\[
 \mathcal L=\langle\mathcal C_{\leq-1/2}[1],\mathcal C_{>1/2}\rangle
             =\mathcal B_-\cap\mathcal B_+.
\]
To see the torsion pair $(\mathcal L,\mathcal M)$ in $\mathcal B_-$,
take $E\in\mathcal B_-$ and let $M_E$ be the quotient of $H^0(E)$
obtained by removing its torsion and its slope factors above $1/2$.
Then $M_E\in\mathcal M$.  The kernel of $E\to M_E$ in
$\mathcal B_-$ has degree-minus-one cohomology $H^{-1}(E)$ and
degree-zero cohomology in $\mathcal C_{>1/2}$, so it lies in
$\mathcal L$.  The slope vanishing and the ordinary cohomological
degrees give $\Hom(\mathcal L,\mathcal M)=0$.
Tilting this torsion pair gives
$\mathcal B_+=\langle\mathcal M[1],\mathcal L\rangle$; in
$\mathcal B_+$ the corresponding torsion pair is
$(\mathcal M[1],\mathcal L)$.  In particular
$W_-[1]\subset W_+\subset W_-$, and
$W_+[1]\subset W_-[1]\subset\mathcal U_-^\theta$.
It suffices to prove
\begin{equation}\label{eq:comparison-hom}
               \Hom(\mathcal T_+^\alpha,\mathcal F_-^\theta)=0.
\end{equation}
Indeed an object $E\in W_-$ has maps to a member of
$\mathcal F_-^\theta$ exactly through $H^0_{\mathcal B_-}(E)$;
vanishing of all such maps puts this cohomology object in
$\mathcal T_-^\theta$.  Equation~\eqref{eq:comparison-hom} therefore
puts $\mathcal T_+^\alpha$ in $\mathcal U_-^\theta$, and the preceding
inclusion handles the negative $\mathcal B_+$ degrees.

Suppose there is a nonzero map $E\to G$ with
$E\in\mathcal T_+^\alpha$ and $G\in\mathcal F_-^\theta$.
Its source decomposes in $\mathcal B_+$ as
$0\to M_E[1]\to E\to E_0\to0$, with $E_0\in\mathcal L$.
Because $\Hom(\mathcal M[1],\mathcal B_-)=0$, the map factors through
$E_0$, which is still in $\mathcal T_+^\alpha$.  The
$(\mathcal L,\mathcal M)$ decomposition of $G$ in $\mathcal B_-$
then factors it through a subobject $G_0\in\mathcal L$ of $G$, since
$\Hom(\mathcal L,\mathcal M)=0$.

Take the kernel $K$ and the nonzero image $Q$ of $E_0\to G_0$ in
$\mathcal B_-$.  Let $K_{\mathcal L}$ be the $\mathcal L$ torsion
part of $K$, and put $U'=E_0/K_{\mathcal L}$ in $\mathcal B_-$.
The quotient closure of $\mathcal L$ gives $Q,U'\in\mathcal L$.
All three terms of
$K_{\mathcal L}\to E_0\to U'$ belong to both hearts, so its triangle
is short exact in $\mathcal B_+$ as well.  Thus $U'$ remains a
quotient of $E$ in $\mathcal B_+$, and we have
\begin{equation}\label{eq:common-heart-sequence}
 0\longrightarrow S\longrightarrow U'\longrightarrow Q\longrightarrow0
 \quad\text{in }\mathcal B_-,\qquad
 S=K/K_{\mathcal L}\in\mathcal M,
 \quad U'\in\mathcal T_+^\alpha,
 \quad 0\ne Q\in\mathcal F_-^\theta.
\end{equation}
We use the sequence~\eqref{eq:common-heart-sequence} only as a short exact
sequence in $\mathcal B_-$; its term $S$
need not belong to $\mathcal B_+$.

We now use the slope envelope to exclude factors close to the two slope
boundaries. The resulting rank bounds will contradict the opposite signs of
the two cut forms on $U'$ and $Q$.

Directly from the three affine functions defining $f$,
\begin{equation}\label{eq:strip-test}
 \max_{1/2\leq s\leq s_0}\bigl(f(s)-s/2+u\bigr)
                 =u-\frac{13857}{50000}.
\end{equation}
The perturbation error on a slope $\pm s$ is at most
$(1+s_0)\delta$, so \eqref{eq:strip-margin} makes the relevant
sheaf value of $N_-$ on $[-s_0,-1/2]$, or of $N_+$ on
$[1/2,s_0]$, strictly negative.

Write $A=H^{-1}(Q)$.  If $A\ne0$, its top slope factor $A_1$
is a saturated subsheaf and $A/A_1\in\mathcal C_{\leq-1/2}$.
Consequently $A_1[1]$ is a subobject of $A[1]$, and then of $Q$, in
$\mathcal B_-$.  If $\mu(A_1)\in[-s_0,-1/2]$, the preceding
negative sheaf value makes $A_1[1]$ high, contradicting
$Q\in\mathcal F_-^\theta$.  Thus $\mu^+(A)<-s_0$.

For $U'$, put $A'=H^{-1}(U')$ and $D'=H^0(U')$.
If $D'$ has positive rank, remove its ordinary torsion subsheaf and
take its bottom positive-rank slope quotient $G'$.  The kernel of
$D'\to G'$ consists of that torsion and the other slope factors,
all of slope $>1/2$; this sequence is therefore exact in
$\mathcal B_+$.  The ordinary cohomology sequence also makes $D'$
a $\mathcal B_+$ quotient of $U'$, so $G'$ is such a quotient.
A slope $\mu(G')\in(1/2,s_0]$ would give
$p_+(G')>0$, $N_+^\alpha(G')<0$, contradicting
$U'\in\mathcal T_+^\alpha$.  Every positive-rank slope of $D'$
therefore exceeds $s_0$.

For any $L\in\mathcal L$, its two ordinary cohomology sheaves give
\begin{equation}\label{eq:overlap-rank-bound}
 c(L)\geq\frac12\bigl(r(H^{-1}(L))+r(H^0(L))\bigr)
                                     \geq\frac12|r(L)|.
\end{equation}
Divisorial torsion contributes nonnegative $c$, and torsion in lower
dimension contributes zero.  The stronger slope conclusions just
proved similarly give
\begin{equation}\label{eq:two-slope-tests}
          c(Q)\geq-s_0r(Q),\qquad c(U')\geq s_0r(U').
\end{equation}
More explicitly, with $D=H^0(Q)$ the estimates are
\[
 \begin{split}
 c(Q)&\geq s_0r(A)+\tfrac12r(D)\geq-s_0r(Q),\\
 c(U')&\geq\tfrac12r(A')+s_0r(D')\geq s_0r(U').
 \end{split}
\]
They include torsion cohomology sheaves: their ranks are zero and
their contributions to $c$ are nonnegative.  Also $d(S)\leq0$, by applying
Lemma~\ref{lem:slope-input} to the ordinary slope factors of $S$.

Set $C=c(Q)+c(U')$.  By \eqref{eq:overlap-rank-bound}, $C\geq0$.
If $C=0$, then $c(Q)=r(Q)=0$, so $p_-(Q)=0$, impossible for a
nonzero member of $\mathcal F_-^\theta$.  Hence $C>0$; and
\eqref{eq:two-slope-tests} together with
\eqref{eq:common-heart-sequence} gives $C\geq s_0r(S)$.
By additivity,
\begin{align*}
 N_-^\theta(Q)-N_+^\alpha(U')
 &=\frac12 C-u r(S)-d(S)
       +\theta_1c(Q)+\theta_0r(Q)-\alpha_1c(U')-\alpha_0r(U')\\
 &\geq \left(\frac12-\frac{u}{s_0}-3\delta\right)C-d(S)>0.
\end{align*}
Here the absolute value of the last four terms is at most $3\delta C$
by \eqref{eq:overlap-rank-bound}; this applies even when a derived
rank is negative.  The strict sign is \eqref{eq:comparison-margin}.
But $N_-^\theta(Q)\leq0\leq N_+^\alpha(U')$, a contradiction.
This proves \eqref{eq:comparison-hom} and the aisle inclusion.
\end{proof}

\begin{lemma}\label{lem:twist-aisle}
For all $\alpha,\theta\in B_\delta$,
\[
 \OO_X\in\mathcal A_-^\theta,\qquad
 \OO_X[2]\in\mathcal A_+^\alpha,\qquad
 T_{\OO_X}(\mathcal U_+^\alpha)\subset\mathcal U_-^\theta.
\]
\end{lemma}

\begin{proof}
The object $\OO_X\in\mathcal B_-$ has the smallest positive value
$p_-=1/2$ and $N_-^\theta=u+\theta_0>0$.
It has no nonzero subobject of $p_-=0$: a $p_-=0$ object is an
extension of a sheaf of dimension at most one by $F[1]$ of slope
$-1/2$, and both have zero Hom to $\OO_X$.
Any nonzero low quotient of $\OO_X$ would have $p_-=1/2$, leaving
a kernel of $p_-=0$; that kernel must vanish, contradicting that
$\OO_X$ is high.  Hence $\OO_X\in\mathcal T_-^\theta$.

Similarly $\OO_X[1]\in\mathcal B_+$ has $p_+=1/2$ and
$N_+^\alpha=-u-\alpha_0<0$.  A $p_+=0$ object has no map to it:
for its slope-$1/2$ part this is
$\Hom(F[1],\OO_X[1])=\Hom(F,\OO_X)=0$, and for its
dimension-at-most-one part $T$ it follows from
\[
 \Ext^1(T,\OO_X)=H^2(X,T)^*=0.
\]
Here we use Serre duality and $K_X\simeq\OO_X$.  Every nonzero
subobject $A$ of $\OO_X[1]$ therefore has $p_+(A)=1/2$, and its
quotient $Q$ has $p_+(Q)=0$.  Hence
$N_+^\alpha(A)=N_+^\alpha(\OO_X[1])-N_+^\alpha(Q)<0$.
This proves $\OO_X[1]\in\mathcal F_+^\alpha$, and thus
$\OO_X[2]\in\mathcal A_+^\alpha$.

For $B\in\mathcal U_+^\alpha$ and $j\geq2$, Serre duality and
the heart orthogonality give
\[
 \Ext^j(\OO_X,B)
       =\Hom(B[j-1],\OO_X[2])^*=0,
\]
since $B[j-1]\in\mathcal U_+^\alpha[1]$.
Lemma~\ref{lem:compare-slope-cuts} puts $B$ in
$\mathcal U_-^\theta$.  In the twist triangle
\[
 B\longrightarrow T_{\OO_X}(B)\longrightarrow
          \operatorname{RHom}(\OO_X,B)\otimes\OO_X[1]\longrightarrow B[1],
\]
the third term is a finite direct sum of copies of
$\OO_X[1-j]$ with $j\leq1$, because bounded complexes of
finite-dimensional vector spaces split into their cohomology.
All these terms lie in $\mathcal U_-^\theta$.  Extension closure
therefore gives the claimed twist inclusion.
\end{proof}

\begin{proposition}[Uniform independent-parameter inclusion]
\label{prop:uniform-aisle}
There is a fixed open square
\begin{equation}\label{eq:fixed-omega}
 \Omega=\{\eta=(\eta_0,\eta_1):|\eta_0|,|\eta_1|<5\cdot10^{-6}\}
\end{equation}
and bounded aisles $U_\eta=\mathcal U_-^\eta$, with hearts
$\mathcal A_\eta=U_\eta\cap(U_\eta[1])^\perp$, such that
\begin{equation}\label{eq:uniform-aisle}
                    \Phi U_\eta\subset U_\theta
                    \qquad(\eta,\theta\in\Omega).
\end{equation}
Moreover
\begin{equation}\label{eq:perturbed-positive-form}
 N_\eta=d+\frac12c+ur+\eta_1c+\eta_0r\geq0
                    \quad\text{on }\mathcal A_\eta,
 \qquad N_0=\frac{\Im Z}{5t_0},
\end{equation}
and no nonzero object of any $\mathcal A_\eta$ has zero numerical
class.  The two parameters in \eqref{eq:uniform-aisle} vary
independently throughout this same square.
\end{proposition}

\begin{proof}
Tensoring by $\OO_X(1)$ takes $\mathcal B_-$ to $\mathcal B_+$
and preserves the $p$ test.  Its effect on the other test is
\[
 N_+^\alpha(E(1))=N_-^\eta(E)
 \quad\text{when}\quad
             \alpha_1=\eta_1,\qquad\alpha_0=\eta_0-\eta_1.
\]
Thus it takes the second torsion pair and aisle to those with this
parameter $\alpha$.  If $\eta\in\Omega$, then $\alpha\in B_\delta$;
also $\theta\in\Omega$ implies $\theta\in B_\delta$.
Lemma~\ref{lem:twist-aisle} now gives
\[
 \Phi U_\eta=T_{\OO_X}\bigl(\mathcal U_+^\alpha\bigr)
                                      \subset U_\theta.
\]
The remaining assertions are Lemma~\ref{lem:second-tilt} and
\eqref{eq:imaginary-charge}.

The parameter shear was used once, in proving this inclusion for
the original family.  In particular, applying any power $\Phi^i$
to \eqref{eq:uniform-aisle} gives
$\Phi^{i+1}U_\eta\subset\Phi^iU_\theta$ with the same
$\eta,\theta\in\Omega$; it performs no further shear and requires
no further shrinkage.  Negative powers are allowed by
Lemma~\ref{lem:periodicity}.
\end{proof}
